\cleardoublepage
%\appendix
%\renewcommand{\thechapter}{B}

%\renewcommand{\thesection}{\Alph{chapter}.\arabic{section}}



\chapter{Boxenverzeichnis}
\label{anhangB}


%\chapter{Boxenverzeichnis}


\thispagestyle{empty}


%\addcontentsline{toc}{chapter}{Anhang B – Boxenverzeichnis}

\section{Einführung: Warum Mathematik \newline Struktur ist}
\begin{NoteBox}[Hinweisboxen]

		\begin{itemize}
	\item \emph{Unabhängige Entdeckungen}\dotfill \pageref{box:entdeckung}
	\item \emph{Experiment als Übersetzer}\dotfill \pageref{box:experiment}
	

	\item \emph{Zahlen als universelle Entdeckung}\dotfill \pageref{box:universelle}



		\end{itemize}
\end{NoteBox}

\section{Zahlen und Wirklichkeit}
\begin{NoteBox}[Hinweisboxen]
	
	\begin{itemize}
		
		\item \emph{Mathematik als Entdeckung}\dotfill \pageref{box:mathentdeckung}
		\item \emph{Zahlen als universelle Struktur}\dotfill \pageref{box:negative}
		
	\end{itemize}
\end{NoteBox}

\begin{HistoryBox}[Historische Box]
	
	\begin{itemize}
		\item \emph{Der Ishango-Knochen}\dotfill \pageref{box:ishango}
		
	\end{itemize}
\end{HistoryBox}
\begin{MathBox}[Mathematische Boxen]
	
	\begin{itemize}
		\item \emph{Widerspruchsbeweis für  $\sqrt{2}$}\dotfill \pageref{box:widerspruch}
			\item \emph{Rechenregeln für negative Zahlen}\dotfill \pageref{box:negative}
	\end{itemize}
\end{MathBox}





\section{Algebra und Gleichungen}
\begin{HistoryBox}[Historische Boxen]
	
	\begin{itemize}
		\item \emph{Von der Praxis zur Theorie}\dotfill \pageref{box:praxis}
			\item \emph{Vom praktischen Problem zur abstrakten Gleichung}\dotfill \pageref{box:abstraktegleichung}
			
	\end{itemize}
\end{HistoryBox}

\begin{NoteBox}[Hinweisboxen]
	
	\begin{itemize}
		\item \emph{Ein gemeinsames Erbe}\dotfill \pageref{box:erben}
		\item \emph{Was Algebra wirklich leistet}\dotfill \pageref{box:allgemeinform}
	\end{itemize}
\end{NoteBox}

\begin{DidacticBox}[Didaktikboxen]
	\begin{itemize}
	\item \emph{Der Ursprung des Wortes „Algorithmus“}\dotfill \pageref{box:ursprung}
	\item \emph{Warum Gleichungen die Realität abbilden}\dotfill \pageref{box:realitaet}
		\item \emph{Warum Regeln die Mathematik tragen}\dotfill \pageref{box:regeln}
\end{itemize}
\end{DidacticBox}


\section{Mathematik als Sprache der Natur}
\begin{HistoryBox}[Historische Boxen]
	
	\begin{itemize}
		\item \emph{Vom Deuten zum Messen}\dotfill \pageref{box:messen}
		\item \emph{Universelle Gesetze statt lokaler Regeln}\dotfill \pageref{box:gesetze}
	\end{itemize}
\end{HistoryBox}

\begin{DidacticBox}[Didaktikboxen]
	\begin{itemize}
		\item \emph{Was eine Differentialgleichung beschreibt}\dotfill \pageref{box:differential}
		\item \emph{Gleiche Mathematik, neue Begriffe}\dotfill \pageref{box:neuebegriffe}
	
	\end{itemize}
\end{DidacticBox}


\section{Abstraktion und Struktur}
\begin{DidacticBox}[Didaktikboxen]
	\begin{itemize}
		\item \emph{Was Symmetrie bedeutet}\dotfill \pageref{box:symmetrie}
	
			\item \emph{Abstraktion und Realität}\dotfill \pageref{box:abstraktion}
	\end{itemize}
\end{DidacticBox}




\section{Zufall -- Eigenschaft der Welt oder Grenze unserer Beschreibung?}


\begin{DidacticBox}[Didaktikboxen]
	\begin{itemize}
		\item \emph{Einzelwurf vs. Gesetz im Ganzen}\dotfill \pageref{box:einzelwurf}
		\item \emph{Statistik ist kein Notbehelf}\dotfill \pageref{box:statistik}
		\item \emph{Was die Quantenmechanik wirklich vorhersagt}\dotfill \pageref{box:vorhersagt}
	\end{itemize}
\end{DidacticBox}

\section{Wie komplexe Zahlen Struktur erzeugen}

\begin{HistoryBox}[Historische Box]
	
	\begin{itemize}
		\item \emph{Historischer Kontext: Julia und Fatou}\dotfill \pageref{box:julia}
	\end{itemize}
\end{HistoryBox}

\begin{NoteBox}[Hinweisboxen]
	
	\begin{itemize}
		\item \emph{Wichtige Unterscheidung: \newline Parameterraum vs.\ Realraum}\dotfill \pageref{box:unterscheidung}
		\item \emph{Kernaussage dieses Kapitels}\dotfill \pageref{box:kernaussage}
	\end{itemize}
\end{NoteBox}

\begin{DidacticBox}[Didaktikboxen]
	\begin{itemize}
		\item \emph{Stabil vs.\ chaotisch: nicht verwechseln}\dotfill \pageref{box:stabil}
		\item \emph{Der häufigste Denkfehler}\dotfill \pageref{box:denkfehler}
		\item \emph{Selbstähnlichkeit heißt nicht: identische Kopie}\dotfill \pageref{box:kopie}
		\item \emph{Optische Ähnlichkeit ist kein Modell}\dotfill \pageref{box:optisch}
	\end{itemize}
\end{DidacticBox}

\begin{MathBox}[Mathematische Boxen]
	
	\begin{itemize}
		\item \emph{Begriffe: Orbit und Fixpunkt}\dotfill \pageref{box:orbit}
		\item \emph{Beschränkt oder unbeschränkt}\dotfill \pageref{box:beschraenkt}
		\item \emph{Definition: Landkarte im Parameterraum}\dotfill \pageref{box:landkarte}
		\item \emph{Definition: Iteration, gefüllte \newline Julia-Menge, Julia-Rand}\dotfill \pageref{box:iteration}
		\item \emph{Box-Counting-Dimension}\dotfill \pageref{box:boxcounting}
	\end{itemize}
\end{MathBox}


\section{Mathematik und Philosophie}

\begin{HistoryBox}[Historische Boxen]
	
	\begin{itemize}
		\item \emph{Platon: Mathematik als Zugang zur Ideenwelt}\dotfill \pageref{box:platon-ideenwelt}
		\item \emph{Aristoteles: Mathematik als Abstraktion}\dotfill \pageref{box:aristoteles-abstraktion}
		\item \emph{Immanuel Kant und die \glqq Kritik der reinen Vernunft\grqq}\dotfill \pageref{box:kant-krv}
	\end{itemize}
\end{HistoryBox}

\begin{DidacticBox}[Didaktikboxen]
	\begin{itemize}
		\item \emph{Platonisches Denken am Quadrat}\dotfill \pageref{box:platon-quadrat}
		\item \emph{Aristoteles im Ingenieurstil}\dotfill \pageref{box:aristoteles-ingenieur}
		\item \emph{Kants Kernidee am Dreieck}\dotfill \pageref{box:kant-dreieck}
		\item \emph{Klarer Satz ohne Mystik}\dotfill \pageref{box:ohne-mystik}
	\end{itemize}
\end{DidacticBox}

\begin{NoteBox}[Hinweisboxen]
	
	\begin{itemize}
		\item \emph{Der Preis der Erklärung}\dotfill \pageref{box:kant-preis}
		\item \emph{Der strukturalistische Kern}\dotfill \pageref{box:strukturalistischer-kern}
		\item \emph{Arbeitsdefinition für dieses Buch}\dotfill \pageref{box:arbeitsdefinition}
	\end{itemize}
\end{NoteBox}

\section{Mathematik im 21.~Jahrhundert}

\begin{DidacticBox}[Didaktikboxen]
	\begin{itemize}
		\item \emph{Die eigentliche Pointe von Gödel}\dotfill \pageref{box:goedel-pointe}
		\item \emph{Computer und Beweis}\dotfill \pageref{box:computer-beweis}
		\item \emph{Experimentelle Mathematik}\dotfill \pageref{box:experimentelle-mathematik}
		\item \emph{KI und mathematische Wahrheit}\dotfill \pageref{box:ki-wahrheit}
		\item \emph{Symmetrie und Erhaltung}\dotfill \pageref{box:symmetrie-erhaltung}
	\end{itemize}
\end{DidacticBox}


\section{Die Mathematik der Natur}

\begin{NoteBox}[Hinweisboxen]
	
	\begin{itemize}
		\item \emph{Die Biene rechnet nicht -- aber die Struktur entsteht}\dotfill \pageref{box:biene-struktur}
		\item \emph{Fibonacci und Pflanzen}\dotfill \pageref{box:fibonacci-pflanzen}
	\end{itemize}
\end{NoteBox}

\begin{DidacticBox}[Didaktikboxen]
	\begin{itemize}
		\item \emph{Natur und mathematische Struktur}\dotfill \pageref{box:natur-struktur}
		\item \emph{Nicht Zufall, sondern funktionale Struktur}\dotfill \pageref{box:verzweigungsnetzwerk}
		\item \emph{Fibonacci nicht mystisch verstehen}\dotfill \pageref{box:fibonacci-mystisch}
		\item \emph{Warum Seifenhäute mathematisch interessant sind}\dotfill \pageref{box:seifenhaut}
		\item \emph{Symmetrie und Variation}\dotfill \pageref{box:symmetrie-variation}
		\item \emph{Ordnung ohne zentralen Bauplan}\dotfill \pageref{box:selbstorganisation}
		\item \emph{Modell und Wirklichkeit}\dotfill \pageref{box:modell-wirklichkeit}
	\end{itemize}
\end{DidacticBox}


\section{Mathematische Hintergründe und Herleitungen}

\begin{DidacticBox}[Didaktikboxen]
	\begin{itemize}
		\item \emph{Warum dieser Beweis so elegant ist}\dotfill \pageref{box:beweis-elegant}
		\item \emph{Warum die Einführung von $i$ notwendig ist}\dotfill \pageref{box:i-notwendig}
		\item \emph{Kernidee der Euler-Formel}\dotfill \pageref{box:euler-kernidee}
		\item \emph{Warum diese „Log-Log“-Formel?}\dotfill \pageref{box:mandelbrot-loglog}
		\item \emph{Merksatz: Nicht Strings im Gehirn -- sondern\newline Flächen-Optimierung}\dotfill \pageref{box:stringtheorie-flaechenoptimierung}
	\end{itemize}
\end{DidacticBox}

\begin{MathBox}[Mathematische Boxen]
	\begin{itemize}
		\item \emph{Grundprinzip: Pixel $\rightarrow$ Parameter $c$ $\rightarrow$ Iteration}\dotfill \pageref{box:mandelbrot-grundprinzip}
		\item \emph{Escape-Radius $R=2$}\dotfill \pageref{box:mandelbrot-escape-radius}
		\item \emph{Vom Draht zum Rohr: der entscheidende Schritt}\dotfill \pageref{box:draht-rohr}
	\end{itemize}
\end{MathBox}

\begin{NoteBox}[Hinweisboxen]
	\begin{itemize}
		\item \emph{Typisches Standardfenster}\dotfill \pageref{box:mandelbrot-standardfenster}
		\item \emph{Wichtig für den Leser}\dotfill \pageref{box:mandelbrot-leserhinweis}
	\end{itemize}
\end{NoteBox}